\documentclass[11pt]{article}

\usepackage[margin=1in]{geometry}
\usepackage{amsmath,amssymb}
\usepackage{booktabs}
\usepackage{graphicx}
\usepackage{hyperref}
\usepackage{enumitem}
\usepackage{xcolor}

\title{Phase 1 Proposal: Lightweight Spike-Based DOA Sector Inference for Embedded Edge AI}
\author{Hossein Molhem\\Kennesaw State University\\EE 6900 Embedded AI}
\date{\today}

\begin{document}
\maketitle

\begin{abstract}
This proposal defines Phase 1 of a staged embedded-AI research and engineering project. The objective is to build a compact, reproducible Python reference prototype for direction-of-arrival (DOA) sector classification using a physics-based uniform linear array (ULA) model, beamspace feature extraction, rate-based spike encoding, and a lightweight leaky-integrate-and-fire (LIF)-inspired inference mechanism. The goal is not to claim a novel DOA estimator in Phase 1. Instead, this phase establishes a defensible baseline that connects array signal processing, sparse temporal encoding, and future embedded deployment on PYNQ or FPGA-oriented platforms.
\end{abstract}

\section{Motivation}

Embedded AI systems must often make decisions under limited compute, memory, energy, and latency budgets. Classical deep-learning models may provide high accuracy, but they are not always appropriate for edge sensing platforms where deterministic operation, reduced arithmetic complexity, and hardware-friendly inference are important. Spike-based or SNN-inspired inference provides a useful design direction because information can be represented temporally and sparsely.

The selected sensing case for Phase 1 is DOA sector classification. This choice is technically meaningful because it is grounded in array signal processing while remaining small enough for a focused Python prototype. A ULA-based DOA model also connects naturally to future radar, MIMO sensing, beamspace processing, and FPGA acceleration work.

\section{Problem Statement}

The Phase 1 problem is to classify the angular sector of a single narrowband source impinging on a uniform linear array. The prototype must compare a conventional non-spiking baseline against a sparse spike-encoded LIF-inspired classifier. The main engineering question is:

\begin{quote}
Can compact beamspace features extracted from a physics-based ULA signal model be converted into sparse spike trains and classified with acceptable accuracy and controllable temporal latency?
\end{quote}

\section{Objectives}

The Phase 1 objectives are:

\begin{enumerate}[label=\arabic*.]
    \item Implement a reproducible ULA signal model for single-source DOA-sector classification.
    \item Extract compact beamspace features from sample covariance matrices.
    \item Train and evaluate a deterministic nearest-centroid non-spiking baseline.
    \item Encode normalized features into sparse spike trains using rate-based encoding.
    \item Implement a lightweight LIF-inspired classifier as a temporal evidence accumulator.
    \item Evaluate accuracy versus SNR and estimate stable-decision latency in time steps.
    \item Sweep LIF time-window and leak parameters to study the accuracy--latency trade-off.
    \item Produce report-ready figures, tables, and a clean project handoff.
\end{enumerate}

\section{Technical Approach}

\subsection{Signal Model}

The dataset will be generated from a narrowband single-source ULA model. For a source at angle $\theta$, the array snapshot matrix is modeled as
\begin{equation}
    X = a(\theta)s + N,
\end{equation}
where $a(\theta)$ is the ULA steering vector, $s$ is the narrowband source waveform, and $N$ is circular complex Gaussian noise. The sample covariance matrix is
\begin{equation}
    R_{xx} = \frac{1}{M}XX^{H},
\end{equation}
where $M$ is the number of snapshots.

\subsection{Feature Extraction}

Beamspace energy features will be computed over a fixed angular grid. These features provide a compact representation of spatial information without implementing a full high-resolution DOA estimator in Phase 1.

\subsection{Baseline Classifier}

A nearest-centroid classifier will be used as the deterministic non-spiking baseline. This baseline is simple, reproducible, and interpretable. Its role is to determine whether the compact beamspace features contain enough information for coarse DOA-sector classification before spike encoding.

\subsection{Spike Encoding and LIF-Inspired Inference}

Normalized beamspace features will be converted into binary spike trains using rate-based encoding. The LIF-inspired classifier will accumulate spike-driven evidence in a leaky membrane state for each class. The predicted class is determined by the final membrane value, while stable-decision latency is estimated by identifying the earliest time step after which the final decision remains unchanged.

\section{Initial Dataset Configuration}

\begin{table}[h]
\centering
\caption{Planned Phase 1 Dataset Configuration}
\begin{tabular}{ll}
\toprule
Parameter & Value \\
\midrule
Array type & Uniform Linear Array \\
Number of elements & $N=8$ \\
Snapshots per sample & $M=128$ \\
Element spacing & $d=\lambda/2$ \\
Source count & 1 \\
DOA range & $-60^{\circ}$ to $+60^{\circ}$ \\
Number of sectors & 4 \\
SNR values & $-10,-5,0,5,10,15,20$ dB \\
Samples per sector per SNR & 100 \\
Total samples & 2800 \\
Train/test split & 70/30, stratified \\
\bottomrule
\end{tabular}
\end{table}

\section{Evaluation Metrics}

Phase 1 will use the following metrics:

\begin{itemize}
    \item overall test accuracy,
    \item accuracy versus SNR,
    \item confusion matrix at a representative SNR,
    \item spike-count statistics,
    \item mean stable-decision latency in time steps,
    \item accuracy--latency trade-off under LIF parameter sweeps.
\end{itemize}

\section{Expected Deliverables}

The expected deliverables are:

\begin{enumerate}[label=\arabic*.]
    \item Python source code for the end-to-end Phase 1 pipeline.
    \item CSV and text result files for reproducibility.
    \item Report-ready figures for accuracy, latency, and trade-off analysis.
    \item A formal Phase 1 report.
    \item A compact IEEE-style paper draft.
    \item A presentation slide deck and transcript.
\end{enumerate}

\section{Expected Contribution}

The expected contribution of Phase 1 is a clean engineering bridge between classical array processing and hardware-aware spike-based inference. The output should be suitable as a course deliverable and as a foundation for later phases involving richer radar/DOA case studies, PYNQ deployment, and Vitis HLS kernel exploration.

\section{References}

\begin{thebibliography}{9}
\bibitem{haykin}
S. Haykin, \emph{Adaptive Filter Theory}. Pearson, 2014.

\bibitem{vanTrees}
H. L. Van Trees, \emph{Optimum Array Processing}. Wiley, 2002.

\bibitem{maass}
W. Maass, ``Networks of spiking neurons: The third generation of neural network models,'' \emph{Neural Networks}, vol. 10, no. 9, pp. 1659--1671, 1997.

\bibitem{roy}
K. Roy, A. Jaiswal, and P. Panda, ``Towards spike-based machine intelligence with neuromorphic computing,'' \emph{Nature}, vol. 575, pp. 607--617, 2019.
\end{thebibliography}

\end{document}
